% Modernised reading — fonds Alexandre Grothendieck, shelfmark 29, batch 7
% (pages 121-140).
%
% This is an INTERPRETATION, not a transcription. It re-reads the pages in
% today's notation and vocabulary, states the hypotheses the manuscript leaves
% implicit, and drops the critical apparatus so that the mathematics reads
% straight through. Everything the brackets and underlines carried moves into a
% footnote. In any dispute of substance the transcription governs: whoever
% cites Grothendieck must cite batch-07.fr.tex.
%
% Scope: the folder was read WHOLE before any of this was written — all eleven
% batches, pages 1-216, in one pass — and one file was then written per batch.
% This batch carries the folder's EARLIEST dated piece — a letter to Serre of
% 18 October 1959 — at page 124, in the middle of the carton. The archive's
% order is not the order of thought, and the reading says so where it matters.
%
% Departures from the page, each footnoted where it occurs:
%   - page 137's minimality condition reads (P,xi) ~ (P',xi') x_G G' where
%     extension of the structure group from G' to G requires x_{G'} G. The
%     transcription flags it; the true statement is written here;
%   - page 136's Proposition states base change and descent of admissibility in
%     an order its own words do not support; the two directions are separated;
%   - page 132's condition (v) carries a struck qualifier the transcription
%     could not read; the condition is given as far as it is legible.
%
% Pass: Opus 5 (claude-opus-5), 2026-08-29 — first pass, unchecked against the
% pages by a human.
\documentclass[11pt,a4paper]{article}
% Le préambule commun à toutes les transcriptions du fonds Grothendieck.
%
% Il tient en une idée : **l'incertitude se compose, elle ne se résout pas.**
% Une transcription qui lisse un mot douteux a détruit la seule chose pour
% laquelle on la faisait — un lecteur doit pouvoir savoir, à chaque endroit,
% ce qui était sur le feuillet et ce qui a été ajouté. Les six macros qui
% suivent sont donc l'appareil critique, et non de la décoration.
%
% Les mêmes commandes sont reconnues par `scripts/render.mjs`, qui produit la
% vue de lecture du site. Une macro ajoutée ici doit l'être aussi là-bas,
% faute de quoi le rendu échouera bruyamment — ce qui est voulu : un rendu
% silencieusement faux est pire qu'un rendu absent.

\ProvidesPackage{grothendieck}[2026/08/08 Transcriptions du fonds Grothendieck]

\usepackage{amsmath,amssymb,amsthm}
\usepackage{mathrsfs}
\usepackage{stmaryrd}
% Les diagrammes commutatifs sont partout dans ce fonds : c'est la langue
% même de la géométrie algébrique de Grothendieck, et une transcription qui
% les rendrait en prose ne transcrirait rien.
\usepackage{tikz-cd}
% Français par défaut — c'est la langue des feuillets — mais l'anglais est
% chargé aussi : la traduction et le résumé sont des documents anglais, et la
% typographie française (espaces avant les deux-points) y serait une faute.
% Ces documents basculent par \selectlanguage{english} en tête de corps.
\usepackage[english,french]{babel}
\usepackage{fontspec}
\usepackage{geometry}
\geometry{a4paper,margin=25mm,marginparwidth=28mm}
\usepackage{marginnote}
\usepackage{xcolor}
% ulem, not soul. `soul`'s \st and \ul are built on a character-by-character
% scan that XeLaTeX's font machinery breaks: the document fails with an
% undefined control sequence rather than degrading. ulem does the same two jobs
% and is Unicode-engine safe. `normalem` stops it hijacking \emph, which the
% transcriptions use for Grothendieck's own emphasis.
\usepackage[normalem]{ulem}
\usepackage{hyperref}
\hypersetup{colorlinks=true,linkcolor=black,urlcolor=black,pdfborder={0 0 0}}

% --- Métadonnées du lot -----------------------------------------------------
% Reprises telles quelles par le script de rendu, qui les affiche en tête de
% la vue de lecture. Elles fixent ce que le fichier prétend couvrir : sans
% elles, une transcription flotte sans point d'attache dans un fonds de seize
% mille feuillets.

\newcommand{\folder}[1]{\gdef\grothFolder{#1}}
\newcommand{\batch}[1]{\gdef\grothBatch{#1}}
\newcommand{\leaves}[2]{\gdef\grothFirst{#1}\gdef\grothLast{#2}}
\newcommand{\foldertitle}[1]{\gdef\grothFoldertitle{#1}}
\gdef\grothFolder{?}\gdef\grothBatch{?}\gdef\grothFirst{?}\gdef\grothLast{?}\gdef\grothFoldertitle{}

% --- Le résumé --------------------------------------------------------------
%
% Il ouvre la lecture modernisée, sous le titre « Résumé », et il est composé
% en petit. Ce n'est pas un ornement : c'est le seul passage du document qui
% ne soit pas des mathématiques, et un lecteur doit pouvoir le reconnaître
% avant de l'avoir lu — puis le sauter s'il connaît déjà le sujet, ou s'y
% arrêter s'il ne le connaît pas. Le filet qui le suit marque la reprise du
% corps.

\newenvironment{resume}
  {\small\setlength{\parskip}{0.45em}}
  {\par\medskip\hrule\medskip}

% --- Mots-clés ---------------------------------------------------------------
%
% En anglais, en clôture du résumé de la lecture modernisée : le vocabulaire
% moderne sous lequel chercher ce que les feuillets construisent. Le manifeste
% du site (`npm run manifest`) les extrait pour étiqueter la cote entière —
% cette ligne est la source unique des tags, il n'y a pas de fichier à côté.
\newcommand{\keywords}[1]{\par\smallskip\noindent
  {\footnotesize\textsc{Keywords} --- \textit{#1}}\par}

% --- L'appareil critique ----------------------------------------------------

% Le numéro de feuillet, en marge. C'est l'ancre qui permet de régler un
% désaccord sur une formule en regardant *un* feuillet plutôt qu'une chemise
% de six cents. Le site s'en sert aussi pour tourner les pages du fac-similé
% quand on fait défiler la transcription.
\newcommand{\leaf}[1]{%
  \marginnote{\footnotesize\color{gray!70}\textsf{#1}}%
  \label{leaf:#1}\ignorespaces}

% Illisible : on ne devine pas. Un mot inventé qui se lit comme les autres est
% le pire résultat possible.
\newcommand{\ill}{\textcolor{red!60!black}{[\,\ldots\,]}}

% Lecture douteuse : le mot est donné, et signalé comme incertain.
\newcommand{\uncertain}[1]{\uline{#1}}

% Ajout de l'éditeur — une lettre manquante, un mot sous-entendu.
\newcommand{\add}[1]{\textcolor{blue!50!black}{[#1]}}

% Biffé par Grothendieck lui-même. Ce qu'il a rayé fait partie du document :
% une rature dit souvent où la pensée a changé de direction.
\newcommand{\struck}[1]{\sout{#1}}

% Note du transcripteur, sur l'état matériel du feuillet.
\newcommand{\note}[1]{\par\smallskip{\small\color{gray!60!black}\textit{[#1]}}\par\smallskip}

% Note marginale de Grothendieck, distincte du corps du texte.
\newcommand{\marginal}[1]{\par\smallskip\noindent\hspace{1em}%
  {\small\color{gray!60!black}#1}\par\smallskip}

% --- Titre ------------------------------------------------------------------

\AtBeginDocument{%
  \begin{center}
    {\large\bfseries Fonds Alexandre Grothendieck --- cote n\textsuperscript{o}~\grothFolder}\\[2pt]
    {\itshape\grothFoldertitle}\\[4pt]
    {\small feuillets \grothFirst--\grothLast\ (lot \grothBatch)}\\[2pt]
    {\footnotesize\color{gray!60!black}Transcription établie d'après le fac-similé numérisé
     par l'Université de Montpellier.\\
     Les lectures incertaines sont soulignées, les illisibles notées [\,\ldots\,],
     les ajouts entre crochets.}
  \end{center}
  \vspace{1em}}

\endinput


\folder{29}
\batch{7}
\pages{121}{140}
\dating{1959-1969}
\watermark{Édition de démonstration\\interprétation personnelle de l'œuvre}
\foldertitle{Groupe fondamental [Autour de SGA 4 et SGA 7] — lecture modernisée}

\begin{document}

\section*{Résumé}

\begin{resume}
Au milieu de ce carton, à la page 124, se trouve la pièce la plus ancienne
qu'il contienne : une lettre à Jean-Pierre Serre datée du 18 octobre 1959. Elle
n'a rien à voir avec la ramification modérée des lots précédents. Elle définit
autre chose, et ce quelque chose a un nom aujourd'hui.

Voici la question qu'elle pose. Sur les nombres complexes, un revêtement d'un
espace est déterminé par une permutation d'un ensemble fini : le groupe
fondamental est un groupe ordinaire, et il agit sur des points. En
caractéristique $p$, il existe des objets qui ressemblent à des revêtements de
degré $p$ mais n'ont qu'\emph{un seul point} : ce sont les schémas en groupes
infinitésimaux, dont l'exemple le plus simple est le noyau du Frobenius. Un
tel objet a une taille --- il est de degré $p$ --- et pourtant aucun point à
permuter. Le groupe fondamental ordinaire, qui ne sait compter que des points,
ne les voit pas.

Grothendieck écrit à Serre qu'il a « réfléchi un petit peu à la partie
\emph{infinitésimale} du groupe fondamental, juste assez pour me convaincre que
ça existe et est raisonnable », et il donne le procédé. Le procédé est celui
qu'il emploiera toute sa vie : ne pas construire l'objet, mais construire le
\emph{foncteur} dont il devra être la solution, et montrer que ce foncteur est
représentable --- ici, pro-représentable, ce qui veut dire représentable par une
suite d'approximations finies. Le foncteur est celui qui, à un groupe fini
$G$, associe l'ensemble des revêtements principaux de groupe $G$ munis d'un
point marqué ; il commute aux produits et aux noyaux, et cela suffit
formellement à produire un système projectif $(G_{i})$ qui joue le rôle du
groupe fondamental. Sur un corps de base, il l'appelle « le groupe
proalgébrique (mais avec partie infinitésimale) fondamental du $k$-schéma
$S$ ».

Puis il le calcule sur un exemple, et le calcul est frappant. Pour une variété
abélienne $X$ sur un corps algébriquement clos,
\[
  \pi_{1}(X/k) \;=\; \varprojlim_{n} {}_{n}X ,
\]
la limite projective des noyaux de la multiplication par $n$ --- \emph{avec
leur partie infinitésimale} --- et ce groupe est, au sens de Cartier, dual du
groupe ind-algébrique limite des ${}_{n}X^{*}$ sur la variété duale. C'est,
mot pour mot, ce qu'on appelle aujourd'hui le \emph{schéma en groupes
fondamental}, et son calcul sur les variétés abéliennes. Le nom et la théorie
publiés sont de Madhav Nori, dix-sept ans plus tard.

La lettre finit sur un calendrier, qu'il faut citer pour ce qu'il est : « Sauf
difficultés imprévues ou enlisement, le multiplodoque devrait être fini d'ici
3 ans, ou 4 au maximum. On pourra commencer à faire de la géométrie
algébrique ! »

Autour de la lettre, ce lot contient la fin du tapuscrit sur les groupes
d'inertie --- qui aboutit à une classe de conjugaison d'homomorphismes bien
définie, un objet par composante du diviseur --- et onze feuillets manuscrits
que la chemise appelle « feuilles anciennes », où la même théorie est reprise à
la main. Ils sont difficiles à lire, et ils démontrent deux fois le même
énoncé : sur un fibré principal \emph{minimal}, un homomorphisme de groupes
structuraux est déterminé par son effet. C'est la clef de la
pro-représentabilité, et c'est de là que tout part.

Les noms modernes à chercher sont \emph{schéma en groupes fondamental},
\emph{groupe fondamental de Nori}, \emph{dualité de Cartier}, \emph{foncteur
pro-représentable}.

\keywords{fundamental group scheme, Nori fundamental group, infinitesimal group scheme, Cartier duality, pro-representable functor}
\end{resume}

\pagerange{121}{122}
\section*{Une classe de conjugaison par branche (pages 121 et 122)}

Le tapuscrit du lot 6 s'achève. Sous les conditions du numéro 5 --- $D =
f^{*}(E)$ régulier, $\operatorname{supp} E$ irréductible et régulier aux points
de $f(\operatorname{supp} D)$, avec la condition de transversalité --- les
images par $\varphi : G \to H$ des sous-groupes d'inertie de $G$ en un point
$x$ de $\operatorname{supp} D$ sont des sous-groupes d'inertie de $H$ relatifs
au point maximal de $\operatorname{supp} E$.

Il en résulte le corollaire pratique : pour toute représentation linéaire de
$G$ dans un module $M$ qui se factorise par $H$, \emph{les images dans
$\operatorname{Aut}(M)$ des sous-groupes d'inertie de $G$ relatifs aux divers
points de $\operatorname{supp} D$ sont conjuguées}. C'est ce qu'on veut savoir
quand on étudie une famille : la monodromie locale est la même, à conjugaison
près, tout le long du diviseur.

\subsection*{L'énoncé le plus précis}

On peut demander mieux. Pour tout $\xi$ au-dessus d'un point géométrique
$\bar{x}$ de $\operatorname{supp} D$, considérons l'homomorphisme
\[
  \underline{Z}_{\ell}(1)(\bar{x}) \longrightarrow G
\]
dont l'image est le sous-groupe d'inertie correspondant. Pour $\xi$ variable et
$\bar{x}$ fixé, il varie par composition avec un automorphisme intérieur de
$G$ : on a donc, pour $\bar{x}$ fixé, une \emph{classe de conjugaison}
d'homomorphismes.

Pour $\bar{x}$ variable dans une composante $D_{i}$ de l'ensemble des points
réguliers de $D$ de caractéristique $\neq \ell$, on passe d'un tel homomorphisme
à un autre en composant avec l'un des isomorphismes
$\underline{Z}_{\ell}(1)(\bar{x}) \to \underline{Z}_{\ell}(1)(\bar{x}')$
associés à une classe de chemins de $\bar{x}$ à $\bar{x}'$. Si cet isomorphisme
ne dépend pas du choix de la classe de chemins --- c'est-à-dire si le faisceau
$\ell$-adique de Tate $\underline{Z}_{\ell}(1)$ est \emph{trivial} sur
$D_{i}$, ce qui se voit déjà sur sa restriction au point générique --- alors on
identifie tous les $\underline{Z}_{\ell}(1)(\bar{x})$ à un même groupe
$\underline{Z}_{\ell}(1)(D_{i})$, et l'on obtient une classe de conjugaison
d'homomorphismes
\[
  \underline{Z}_{\ell}(1)(D_{i}) \longrightarrow G
\]
bien déterminée par la seule composante $D_{i}$.

Si de plus $\underline{Z}_{\ell}(1)$ est trivial sur l'ouvert $X_{\ell}$ des
points de $X$ de caractéristique $\neq \ell$ et que celui-ci est connexe, on
identifie tout à un unique $\underline{Z}_{\ell}(1)$, et l'on obtient, pour
chaque $i$, une classe de conjugaison d'homomorphismes
$\underline{Z}_{\ell}(1) \to G$. \emph{Ces classes peuvent être distinctes pour
des $i$ distincts} --- et c'est bien pour cela qu'on les a construites. La
formation en est fonctorielle relativement à un morphisme $f : X \to Y$ du type
étudié aux numéros 2 et 4, « ce que le lecteur explicitera lui-même s'il en a
besoin ».

Une note en marge de la page 121, à ne pas laisser passer : « Remords au
n\textsuperscript{o} 3, 4 ». Elle porte sur les deux numéros où la classe
d'inertie est construite, et le carton n'en dit pas davantage.

\pagerange{124}{126}
\section*{Lettre à Serre, 18 octobre 1959 : le groupe fondamental infinitésimal (pages 124 à 126)}

\subsection*{Le contexte, en trois lignes}

La lettre s'ouvre sur Tate et les courbes elliptiques, dont Grothendieck dit
n'avoir « pas du tout compris pourquoi ses résultats suggèreraient l'existence
d'une telle définition » --- il s'agit d'une définition globale des variétés
analytiques sur un corps valué complet --- et reste sceptique. Puis, sans
transition :

\begin{quote}
J'ai réfléchi un petit peu à la partie « infinitésimale » du groupe
fondamental, juste assez pour me convaincre que ça existe et est raisonnable.
\end{quote}

\subsection*{Le dispositif}

On se donne un schéma connexe $S$ (par exemple un schéma algébrique sur un
corps $k$), une catégorie auxiliaire $\mathcal{C}$ (par exemple les schémas
algébriques finis sur $k$) où les produits fibrés existent, et une catégorie
$\mathcal{G}$ dont les objets sont des groupes de $\mathcal{C}$ (par exemple
les groupes algébriques finis sur $k$), soumise à :

\begin{enumerate}
  \item[(i)] $\mathcal{G}$ est stable par produits, et par noyaux de couples de
    morphismes --- le sous-groupe maximum où deux morphismes coïncident, qui
    existe puisqu'il s'exprime par un produit fibré ;
  \item[(ii)] les images existent, sont des quotients, et sont dans
    $\mathcal{G}$ ;
  \item[(iii)] toute suite décroissante de sous-groupes appartenant à
    $\mathcal{G}$ d'un objet de $\mathcal{G}$ est stationnaire.
\end{enumerate}

Et un foncteur covariant $F$ de $\mathcal{G}$ dans les schémas en groupes sur
$S$ (par exemple $G \mapsto S \times_{k} G$), soumis à :

\begin{enumerate}
  \item[(iv)] $F$ est \emph{exact} : il commute aux produits, aux noyaux de
    couples, aux images ;
  \item[(v)] tout $H$ de la forme $F(G)$ est « spécial », c'est-à-dire admet
    une suite exacte de schémas en groupes finis et plats sur $S$
    \[
      e \longrightarrow H_{\mathrm{inf}} \longrightarrow H \longrightarrow
      H_{\mathrm{s\acute{e}p}} \longrightarrow e ,
    \]
    où $H_{\mathrm{inf}}$ est purement infinitésimal --- la projection sur $S$
    est géométriquement injective --- et $H_{\mathrm{s\acute{e}p}}$ non ramifié
    sur $S$ ;
  \item[(vi)] $S$ est réduit.
\end{enumerate}

Il ajoute, et c'est important pour lire ce qui suit : « Les conditions (v) et
(vi) ont une sale gueule, et sont essentiellement provisoires. » Elles ne
servent qu'à un lemme.

\textbf{Lemme.} Soient $H$ comme en (v), $S$ comme en (vi), $P$ et $Q$ deux
fibrés principaux homogènes sous $H$, et $u$, $v$ deux isomorphismes de $P$
dans $Q$ transformant un point marqué donné en un autre donné. Alors $u = v$.

\emph{Démonstration.} En tordant $H$, on se ramène au cas où $P$ est trivial ;
alors il suffit d'observer qu'une section de $Q$ est un isomorphisme de $S$ sur
une composante connexe de $Q_{\mathrm{r\acute{e}d}}$.

Le lemme est ce qui fait tenir toute la construction : \emph{un fibré
principal pointé n'a pas d'automorphisme pointé}. Sans cela, la classification
n'est pas un ensemble mais un groupoïde, et le foncteur ci-dessous ne serait
pas à valeurs dans les ensembles.

\subsection*{Le foncteur, et ce qu'il vérifie}

Soit $a$ un point marqué de $S$ --- une extension algébriquement close du corps
résiduel d'un $s \in S$. Pour $G \in \mathcal{G}$, on note
\[
  Z(S, a\,;\, G) \;=\; Z(G)
\]
l'ensemble des classes d'isomorphie de fibrés principaux homogènes sur $S$, de
groupe $F(G)$, munis d'un point marqué au-dessus de $a$. C'est un foncteur
covariant de $\mathcal{G}$ dans les ensembles, et il vérifie :

\begin{enumerate}
  \item[1)] il commute aux produits, puisque $F$ y commute ;
  \item[2)] il commute aux noyaux de couples : si
    $G'' \to G \rightrightarrows G'$ est exacte dans $\mathcal{G}$, alors
    $Z(G'') \to Z(G) \rightrightarrows Z(G')$ l'est, c'est-à-dire que $Z(G'')$
    s'identifie à l'ensemble des éléments de $Z(G)$ ayant même image par les
    deux flèches. Cela résulte de l'exactitude de $F$ et du lemme.
\end{enumerate}

\subsection*{La pro-représentabilité}

De ces deux propriétés il résulte \emph{formellement} qu'il existe un système
projectif filtrant $(G_{i})_{i \in I}$ dans $\mathcal{G}$, à morphismes de
transition épimorphiques, essentiellement unique, avec un isomorphisme
fonctoriel
\[
  Z(G) \;=\; \varinjlim_{i} \operatorname{Hom}_{\mathcal{G}}(G_{i}, G) .
\]

C'est ce système projectif --- pris modulo l'équivalence qui, intuitivement,
revient à passer à la limite projective des $G_{i}$ --- qui \emph{est} le
groupe fondamental. On le note $\pi_{1}^{\mathcal{G}}(S, a)$ ; et dans le cas
d'un corps de base $k$ avec $\mathcal{G}$ la catégorie des groupes algébriques
finis sur $k$, on écrira $\pi_{1}(S/k, a)$ : « c'est le groupe proalgébrique
(mais avec partie infinitésimale) fondamental du $k$-schéma $S$ ».

Il faut voir la forme du geste, parce qu'il revient dix fois dans le carton :
\emph{on ne construit pas le groupe, on construit le foncteur qu'il devrait
représenter, et on montre que ce foncteur est pro-représentable.} Les pages 164
à 185 feront la même chose pour des groupes structuraux infinis, la page 104
pour $\pi_{2}$, et la page 208 axiomatisera la catégorie où c'est
possible.\footnote{L'existence du système projectif est démontrée à la page 127
par les \emph{couples minimaux}, et c'est le seul endroit où l'argument est
écrit. Un couple $(G, z)$, $z \in Z(G)$, est dit \emph{minimal} si $z$ ne
provient d'aucun $z'$ pour un sous-groupe $G' \subsetneq G$ ; il est
\emph{majoré} par $(G', z')$ s'il existe $u : G' \to G$ avec $z = Z(u)(z')$.
La condition (iii) assure que tout couple est majoré par un couple minimal, et
la propriété 2) que le $u$ est alors \emph{unique}. Les couples minimaux
forment donc un système projectif, filtrant parce que $(G,z)$ et $(G',z')$ sont
tous deux majorés par $(G\times G', (z,z'))$. C'est exactement l'argument que
les feuillets 136 à 140 refont à la main, et deux fois.}

\subsection*{Le dévissage sur un corps parfait}

Lorsque $k$ est parfait, tout groupe algébrique fini se décompose en partie
infinitésimale et partie réduite, et le groupe fondamental est le produit
semi-direct de sa \emph{partie séparable} --- un groupe profini ordinaire sur
lequel $\operatorname{Gal}(\bar{k}/k)$ opère --- par sa \emph{partie
infinitésimale}, laquelle ne dépend d'ailleurs plus du point base.

Deux remarques qu'il faut lire ensemble. La première : la partie séparable est
« plus grande » que le groupe fondamental ordinaire, parce qu'elle classifie
les revêtements principaux dont le groupe structural est un groupe fini
\emph{et séparable sur $k$} --- c'est-à-dire un groupe fini ordinaire sur
lequel on fait opérer $\operatorname{Gal}(\bar{k}/k)$ de façon non
nécessairement triviale. La seconde : « J'avoue que si $k$ n'est pas
algébriquement clos, le groupe fondamental ci-dessus n'est sans doute pas le
bon » ; il faudrait prendre celui de $S \times_{k} \bar{k}$, muni de sa donnée
de descente, donc défini en réalité sur $k$.\footnote{Un paragraphe entier est
biffé à cet endroit, sur les points rationnels du groupe fondamental de
$\operatorname{Spec}(k)$ et le centre de $\operatorname{Gal}(\bar{k}/k)$. La
transcription le conserve, comme elle doit ; il n'est pas repris ici, l'auteur
l'ayant retiré.}

\pagerange{127}{129}
\section*{Ce que la lettre conjecture, et ce qu'elle calcule (pages 127 à 129)}

\subsection*{La suite exacte d'homotopie, à trouver}

« Je ne sais encore comment devrait être formulée la suite exacte d'homotopie ;
pour bien faire, l'inclusion dans le groupe fondamental d'une partie
infinitésimale devrait donner une théorie satisfaisante du comportement du
groupe fondamental par spécialisation. »

L'espoir est le suivant. Soit $f : X \to Y$ propre et séparable à fibres
absolument connexes --- c'est-à-dire $f_{*}(\mathcal{O}_{X}) = \mathcal{O}_{Y}$
--- muni d'une section qui détermine des points-base sur les fibres. Alors les
groupes fondamentaux \emph{totaux} des fibres devraient former sur $Y$ un
pro-schéma en groupes filtrant : il existerait un système projectif
$\pi_{1}(X/Y, s) = (G_{i})_{i \in I}$ de schémas en groupes « spéciaux » sur
$Y$, à morphismes de transition épimorphiques, dont les groupes fondamentaux
des fibres se déduiraient par simple spécialisation.

Et la conjecture proprement dite : $\pi_{1}(X/Y)$, $\pi_{1}(X)$ et l'image
inverse par $f$ de $\pi_{1}(Y)$ devraient être reliés par une suite exacte.
Il ajoute --- et c'est la seule fois du carton où il rêve à voix haute ---
« Il y a de la joie en perspective pour les groupes d'homotopie supérieurs
\dots ».

Une conséquence immédiate est notée : pour des schémas complets $X$, $Y$ sur
$k$ algébriquement clos,
\[
  \pi_{1}(X \times Y/k) \;=\; \pi_{1}(X/k) \times \pi_{1}(Y/k) .
\]
C'est l'énoncé que le lot 4, page 62, démontrera --- mais pour le quotient
premier à la caractéristique, non pour le groupe total. La différence est
exactement la partie infinitésimale, et c'est elle qui rend l'énoncé de 1959
plus fort que celui de 1967.

\subsection*{Le calcul sur une variété abélienne}

C'est le morceau de la lettre.

Utilisant les raisonnements de Serre, on conclut que si $X$ est une variété
abélienne sur $k$, alors $\pi_{1}(X/k)$ est \emph{abélien}, et qu'un revêtement
principal minimal $X'$ de $X$ est une variété abélienne isogène à $X$. On en
déduit
\[
  \pi_{1}(X/k) \;=\; \varprojlim_{n} {}_{n}X ,
\]
où ${}_{n}X$ est le noyau de la multiplication par $n$ --- \emph{avec sa partie
infinitésimale, bien sûr} --- l'homomorphisme de ${}_{mn}X$ dans ${}_{n}X$
étant induit par la multiplication par $m$.

Et, par la dualité de Cartier : ce groupe fondamental est dual du groupe
ind-algébrique
\[
  \varinjlim_{n} {}_{n}X^{*} ,
\]
où $X^{*}$ est la variété duale.

Prenant la $p$-composante, « on trouve ce qui, pour le nombre premier $p$, doit
jouer le rôle du module de Weil » --- c'est-à-dire l'analogue en
caractéristique $p$ du module de Tate. Il ajoute qu'il est hors de doute qu'on
peut associer fonctoriellement à un groupe proalgébrique infinitésimal abélien
un module sur l'anneau de Witt, et que ce module se décompose en les trois
morceaux correspondant aux trois types principaux de groupes algébriques
abéliens ; ce qui donnerait au théorème de Serre « une formulation plus
naturelle », avec une démonstration uniforme sans distinction entre
$\ell \neq p$ et $\ell = p$, et montrerait que « ta somme abracadabrante se
comporte bien quand on fait varier $X$ dans une famille ».

\subsection*{Ce que c'est, aujourd'hui}

L'objet construit dans cette lettre --- le pro-schéma en groupes finis qui
classifie les torseurs pointés sous les schémas en groupes finis, avec sa
partie infinitésimale --- est ce qu'on appelle aujourd'hui le \emph{schéma en
groupes fondamental}. La construction publiée, par la voie tannakienne des
fibrés essentiellement finis, est due à Madhav Nori (thèse de 1976, article de
1982), qui démontre en particulier que pour une variété abélienne
$\pi^{N}(X) = \varprojlim_{n} {}_{n}X$, et la dualité de Cartier avec
$\varinjlim {}_{n}X^{*}$.\footnote{Les deux constructions ne sont pas la même :
celle de Nori est tannakienne et suppose $X$ propre réduit connexe avec un
point rationnel ; celle de la lettre est par les torseurs et repose sur les
conditions (v) et (vi), dont Grothendieck écrit lui-même qu'elles « ont une
sale gueule, et sont essentiellement provisoires ». La lettre est une esquisse
avec des hypothèses avouées comme provisoires, non une théorie. Ce que la
présente note constate est que l'objet, la méthode de pro-représentabilité, le
dévissage semi-direct sur un corps parfait, le calcul sur une variété abélienne
et la dualité de Cartier sont dans une lettre du 18 octobre 1959 restée
inédite. Elle ne constate rien de plus : aucune filiation n'est ici affirmée,
et la question de savoir ce que la littérature publiée établit à ce sujet n'est
pas de la compétence de cette édition.}

\subsection*{Le calendrier}

La lettre se termine ainsi, et il faut la lire jusqu'au bout :

\begin{quote}
J'espère arriver dans l'année prochaine à une théorie satisfaisante du groupe
fondamental, et achever la rédaction des Chapitres IV, V, VI, VII (ce dernier
étant le groupe fondamental), en même temps que des catégories. Dans deux ans
résidus, dualité, intersections, Chern, Riemann-Roch. Dans trois ans
cohomologie de Weil, et un peu d'homotopie si Dieu veut. Et entre-temps, je ne
sais quand, le grand théorème d'existence avec Picard etc, un peu de courbes
algébriques, les schémas abéliens. Sauf difficultés imprévues ou enlisement, le
multiplodoque devrait être fini d'ici 3 ans, ou 4 au maximum. On pourra
commencer à faire de la géométrie algébrique !
\end{quote}

\pagerange{130}{134}
\section*{Les « feuilles anciennes » : représentabilité et recollement (pages 130 à 134)}

Onze feuillets manuscrits que la chemise appelle « \emph{Vrai} groupe
fondamental [feuilles anciennes] » reprennent la même théorie de sa main. Ils
sont d'une écriture pâle et lourdement raturée, et la transcription y laisse
beaucoup d'illisible ; ce qui suit est ce qui s'énonce.

\subsection*{Le foncteur des fibrés pointés}

Soit $f : X \to S$ propre et plat, muni d'une section $\sigma$, avec
$f_{*}(\mathcal{O}_{X}) = \mathcal{O}_{S}$ \emph{universellement}. Soit $G$ un
groupe fini et étale sur $S$ --- le cas important étant $G = S \times \Gamma$
pour $\Gamma$ un groupe fini ordinaire. On note $F_{G}(T)$ l'ensemble des classes d'isomorphie de fibrés principaux
homogènes sur $X_{T}$, de groupe $G_{X_{T}}$, munis d'une section le long de
$\sigma_{T}$ --- c'est-à-dire d'un isomorphisme
$\sigma_{T}^{*}(X') \simeq G_{T}$.

\textbf{Question.} Le foncteur $T \mapsto F_{G}(T)$ est-il représentable ? Si
oui, le schéma qui le représente est \emph{étale} sur $S$. On note que
$F_{G}$ est compatible aux limites et à la descente fidèlement plate
quasi-compacte.

C'est le même foncteur que le $Z(S,a;G)$ de la lettre, rendu relatif : la
lettre le prenait sur un $S$ fixe et cherchait à le pro-représenter ; ici on le
prend en famille et l'on cherche à le représenter. La différence n'est pas de
degré : la seconde question est celle de la \emph{constructibilité en famille}
que la page 44 appelait de ses vœux.

\subsection*{Un critère de représentabilité}

\textbf{Critère.} Soit $S$ un préschéma localement noethérien et
$F : (\mathrm{Sch})_{/S} \to (\mathrm{Ens})$ un foncteur vérifiant :

\begin{enumerate}
  \item[(i)] $F$ est compatible aux limites inductives d'anneaux ;
  \item[(ii)] $F$ est compatible à la descente fidèlement plate quasi-compacte
    et à la localisation ;
  \item[(iii)] si $A$ est un anneau local noethérien \emph{complet} sur $S$ de
    corps résiduel $k$, alors $F(A) \simeq F(k)$ ;
  \item[(iv)] $F$ est séparé ;
  \item[(v)] pour tout corps $k$ sur $S$, $F(k)$ est fini, et
    $F(k) \to F(k')$ est bijectif pour $k'$ une extension algébrique séparable
    de $k$.\footnote{Le qualificatif de $k'$ porte sur la page une mention
    biffée que la transcription ne lit pas ; « alg.\ sép.\ » et « alg.\ clos »
    y voisinent. La condition est donnée ici telle que la page la laisse lire,
    et elle n'est pas complétée par conjecture.}
\end{enumerate}

Alors $F$ est représentable par un schéma \emph{étale, séparé et quasi-fini}
sur $S$.

C'est la condition (iii) qui porte le poids : elle dit que le foncteur ne voit
que la fibre spéciale d'un anneau local complet, ce qui est la version
foncteur du fait qu'un revêtement étale d'un local complet est déterminé par sa
fibre spéciale. La démonstration est esquissée puis s'arrête au milieu d'une
phrase.

\subsection*{Le lemme de recollement}

\textbf{Lemme.} Soit $S$ localement noethérien, et pour tout $s \in S$ un
préschéma $X(s)$ quasi-fini étale séparé sur
$\operatorname{Spec}(\mathcal{O}_{s})$, avec, pour $s \in \overline{\{t\}}$,
des morphismes $\varphi_{st} : X(t) \to X(s)$ au-dessus de
$\operatorname{Spec}(\mathcal{O}_{t}) \to \operatorname{Spec}(\mathcal{O}_{s})$
vérifiant $\varphi_{st}\varphi_{tu} = \varphi_{su}$. Alors il existe un
préschéma $X$ quasi-fini séparé étale sur $S$ et des isomorphismes
\[
  \varphi_{s} : X \times_{S} \operatorname{Spec}(\mathcal{O}_{s})
  \xrightarrow{\ \sim\ } X(s)
\]
compatibles, et $X$ est unique à isomorphisme unique près.

\begin{tikzcd}
X(s) \arrow[d] & X(t) \arrow[l, "\varphi_{st}"'] \arrow[d] \\
\operatorname{Spec}(\mathcal{O}_{s}) & \operatorname{Spec}(\mathcal{O}_{t}) \arrow[l]
\end{tikzcd}

L'unicité résulte d'un énoncé plus général de pleine fidélité sur les
préschémas de type fini sur $S$. Sur l'existence, Grothendieck écrit : « Je ne
suis pas [sûr] de mon raisonnement » --- et la marge propose la stratégie qui
convient : montrer que pour tout $s$, $X(s)$ se réalise sur un voisinage ouvert
$U$ de $s$, puis recoller.

\pagerange{136}{140}
\section*{Groupes admissibles, couples minimaux, et la proposition 4 (pages 136 à 140)}

\subsection*{Les groupes admissibles}

Soit $S$ un schéma \emph{réduit connexe}. Appelons \emph{admissible} un schéma
en groupes $G$, fini et plat sur $S$, admettant une suite exacte
\[
  e \longrightarrow U \longrightarrow G \longrightarrow \Gamma \longrightarrow
  e
\]
avec $U \to S$ plat, fini et géométriquement injectif, et $\Gamma \to S$ plat,
fini et séparable. Alors le sous-groupe $U$ est \emph{déterminé de façon
unique} --- avec, en marge, la question : « pour cette question, l'hypothèse
« $S$ réduit » est-elle nécessaire ? ».

C'est la condition (v) de la lettre, débarrassée du foncteur $F$ : un groupe
est admissible s'il se dévisse en une partie infinitésimale et une partie
étale.

\textbf{Proposition 1.} Soit $S' \to S$ un morphisme et $G$ un $S$-groupe fini.
Si $G$ est admissible, $G \times_{S} S'$ l'est --- c'est trivial, la suite
exacte se changeant de base. Réciproquement, si $G \times_{S} S'$ est admissible
et si $S' \to S$ est fidèlement plat, alors $G$ est admissible, par descente et
unicité de $U$.\footnote{La page énonce les deux directions dans l'ordre
inverse : « Si $G \times_{S} S'$ est admissible, $G$ aussi ; la réciproque est
vraie si $S' \to S$ est fidèlement plat », avec « Direct est trivial » en
dessous. C'est le changement de base qui est trivial et la descente qui demande
la platitude fidèle ; les deux directions sont remises ici dans l'ordre que
l'argument impose. Ce qui est écrit à la suite --- « Réciproque résulte du
théorème de \dots\ et unicité » --- confirme la lecture : c'est bien l'unicité
de $U$ qui sert à descendre.}

\textbf{Corollaire.} Si $G$ est admissible et $P$ un fibré principal sous $G$,
alors $\operatorname{Aut}(P)$, comme $S$-groupe, est admissible.

\textbf{Proposition 2.} Si $(P, \xi)$ est un fibré principal \emph{pointé} sous
un $S$-groupe admissible $G$, alors $\operatorname{Aut}(P, \xi)$ est réduit à
l'identité.

\textbf{Proposition 3.} Deux sections d'un préschéma réduit connexe $P$ sur $S$
qui coïncident en un point sont identiques. Plus précisément : si $s$ est une
section de $P/S$, il existe un sous-préschéma fermé $Q$ de $P$ tel que
$Q_{\mathrm{r\acute{e}d}} \to S$ soit un isomorphisme, d'isomorphisme
réciproque induit par $s$.

Les propositions 2 et 3 sont le lemme de la lettre, découpé en deux : la
seconde est le fait géométrique, la première la conséquence dont on se sert.

\subsection*{Minimalité}

Un couple $(P, \xi)$, relatif à un $G$ admissible, est dit \emph{minimal} s'il
n'existe pas de sous-groupe admissible $G' \subset G$ et de $(P', \xi')$
relatif à $G'$ tels que
\[
  (P, \xi) \;\simeq\; (P', \xi') \times_{G'} G .
\]
Autrement dit : $(P,\xi)$ n'est pas obtenu par \emph{extension du groupe
structural} depuis un sous-groupe strict.\footnote{La page écrit
$(P',\xi') \times_{G} G'$, avec les deux groupes échangés. L'extension du
groupe structural d'un $G'$-fibré à un $G$-fibré est le produit contracté
$P' \times_{G'} G$ : c'est le groupe \emph{du fibré} qui va en indice. La
transcription signale la lecture et la rapproche du $(P,\xi) \times_{G}
(G' \times G')$ des pages 138 et 140, où l'indice est bien le groupe du fibré
--- ce qui confirme que la page 137 est un lapsus et non une autre convention.}

On pose alors
\[
  \mathfrak{Z}(S, a\,;\, G) \;=\; \mathfrak{Z}(G) \;=\;
  \bigl\{\text{classes d'isomorphie de $G$-fibrés principaux pointés au-dessus
  de } a\bigr\} ,
\]
foncteur covariant en le $S$-groupe admissible $G$. \emph{C'est le $Z(S,a;G)$
de la lettre}, écrit d'une autre lettre, et c'est aussi le
$\pi^{1}(S,\xi\,;\,G)$ des pages 164 et suivantes : trois notations, un seul
foncteur, et c'est la lecture du carton entier qui permet de le dire.

\subsection*{La proposition 4, démontrée deux fois}

\textbf{Proposition 4.} Soit $\alpha \in \mathfrak{Z}(G)$ \emph{minimal}. Alors
\[
  u \longmapsto u_{*}(\alpha) , \qquad
  \operatorname{Hom}^{\mathrm{gr}}_{S}(G, G') \longrightarrow \mathfrak{Z}(G')
\]
est \emph{injective}.

\emph{Démonstration.} Soient $u, v : G \to G'$ avec $u_{*}(\alpha) =
v_{*}(\alpha)$. Soit $(P, \xi)$ de classe $\alpha$ et $(P', \xi')$ de classe
$u_{*}(\alpha) = v_{*}(\alpha)$. Le couple $(u,v) : G \to G' \times G'$ donne
\[
  (Q, \eta) \;=\; (P, \xi) \times_{G} (G' \times G')
  \;=\; \bigl[(P,\xi) \times_{u} G'\bigr] \times
        \bigl[(P,\xi) \times_{v} G'\bigr] ,
\]
et par hypothèse il existe un unique homomorphisme $(P', \xi') \to (Q, \eta)$
compatible avec la diagonale $G' \to G' \times G'$.

\begin{tikzcd}
G \arrow[r, "{(u,v)}"] & G' \times G' \\
H \arrow[u] \arrow[r] & G' \arrow[u]
\end{tikzcd}
\qquad
\begin{tikzcd}
(P, \xi) \arrow[r] & (Q, \eta) \\
(R, \lambda) \arrow[u] \arrow[r] & (P', \xi') \arrow[u]
\end{tikzcd}

Soit $(R, \lambda)$ le produit fibré de $(P, \xi)$ et $(P', \xi')$ au-dessus de
$(Q, \eta)$, et $H$ le noyau du couple $G \rightrightarrows G'$ ---
c'est-à-dire le sous-groupe où $u$ et $v$ coïncident. Alors $(R, \lambda)$ est
un fibré principal \emph{de groupe $H$}, et la minimalité de $(P, \xi)$ force
$H = G$, donc $u = v$.

C'est l'énoncé sur lequel toute la pro-représentabilité repose : c'est lui qui
rend \emph{unique} le morphisme entre deux couples minimaux, donc qui fait des
couples minimaux un système projectif et non un groupoïde. La lettre l'utilise
page 127 sans le démontrer ; ces feuillets le démontrent, et deux fois --- la
page 139 porte une première tentative qui s'arrête sur la phrase « Soit
$(R,\lambda)$ le produit fibré de $(P,\xi)$ et $(P',\xi')$ sur $(Q,\eta)$ », et
la page 140 la reprend depuis le début et la conduit à
terme.\footnote{L'abandon de la page 139 est conservé dans la transcription
pour cette raison : il marque l'endroit où l'argument a été refondé. La reprise
de la page 140 change le dispositif --- elle introduit d'emblée la catégorie
$\mathcal{C}$ de $S$-groupes finis et plats stable par produits, noyaux de
couples et images, c'est-à-dire les conditions (i) et (ii) de la lettre --- et
c'est ce cadre-là qui permet de conclure.}

La fin de la démonstration --- l'extension de base $(S_{1}, a_{1}) \to (S, a)$
et l'identification de $P_{1}$ à $G'_{1}$ --- est aux pages 141 et 142, dans le
lot 8, et c'est de là que sort le pro-objet
$\pi_{1}^{\mathcal{C}}(S, \xi)$.

\end{document}
